\documentclass[12pt]{article}
\usepackage[utf8]{inputenc}
\usepackage[T1]{fontenc}
\usepackage{amsmath,amssymb}
\usepackage{geometry}
\geometry{margin=2.5cm}

\begin{document}

\begin{center}
\fbox{\textbf{Theorem (Cauchy)}}
\end{center}

\bigskip
\noindent If $G$ is a finite abelian group and $p\mid$ ord $G$, $p$ prime

\noindent Then $\exists\, a\in G$, ord $a=p$.

\bigskip
\noindent\underline{Proof}:

\noindent Let $n=$ ord $G$, $n\geq p$.

\noindent Induction on $n$.

\bigskip
\noindent $n=p=$ ord $G$.

\noindent $p$ prime $\Rightarrow p>1 \Rightarrow \exists\, a\neq e$. \ ord $a\mid$ ord $G=p$.

\noindent $a\neq e \Rightarrow$ ord $a>1$

\noindent $p$ prime

\noindent $\Rightarrow$ ord $a=p$.

\bigskip
\noindent\big(ord $G=p$ prime $\Rightarrow (\forall)\,a\in G\setminus\{e\}$ we have ord$(a)=p$\big).

\bigskip
\noindent suppose $n>p$, $p\mid n$, and that the statement holds for groups of order $<n$.

\bigskip
\noindent\underline{Case I}: \ $\exists\, a\in G$ such that $p\mid$ ord$(a)$. \ $\Rightarrow$ ord$(a)=p\cdot l$.\footnote{the exact symbol used for this exponent is unclear in the original -- by the mathematics it must be $l=\text{ord}(a)/p$.}

\noindent And $b=a^l \Rightarrow$ ord $b=p$.

\bigskip
\noindent\underline{Case II}: \ $\nexists\, a\in G$ such that $p\mid$ ord$(a)$.

\bigskip
\noindent Let $a\in G$, $a\neq e$. \ suppose ord$(a)=m \Rightarrow p\nmid m$

\noindent Let $H=\{e,\ldots,a^{m-1}\}\trianglelefteq G$ \ (since $G$ is abelian).

\noindent Let $\dfrac{G}{H}$. \ $\Rightarrow \left|\dfrac{G}{H}\right| = [G:H] < n$ \ (by Lagrange)

\bigskip
\noindent ord $G = n = m\cdot[G:H]$.

\[
p\mid n = m\cdot[G:H]
\]
\[
p\nmid m
\]
\[
\Rightarrow p\mid [G:H] = \left|\dfrac{G}{H}\right| < n.
\]

\bigskip
\noindent\underline{induction hyp.} $\Rightarrow \exists\, \hat b\in \dfrac{G}{H}$ such that ord $\hat b=p$; \ $b\in G$, $G$ finite

\[
\Rightarrow \text{ord } b<\infty.
\]
\[
b^t=e
\]
\[
\hat e=\widehat{b^t}=\hat b^t \Rightarrow p\mid t.
\]

\noindent\textit{using the hypothesis from Case II.}\footnote{the mark at the end of the page looks more like a ``q.e.d.''-style marker than a contradiction: $p\mid t=\text{ord}(b)$ is exactly the fact that closes Case II, reducing it to the technique already established in Case I (the element $b^{t/p}$ would have order $p$).}

\end{document}
